\section{Cái chương này đóng được, và cái nó mở ra}
\label{sec:ch08-con-lai-gi}

\index{chi phí!bốn món nợ chương 7 để lại}%
Bốn món nợ chương~\ref{ch:transformer} để lại thì ba trả xong và một trả một
nửa.

Ô \(O(n^2 \cdot d)\) giờ có hằng số. Phần bậc hai vượt nửa một tầng ở
\(n = 2d + \dff\), tức 3072 ở cấu hình của bài, và ở ngưỡng \(n = d\) mà bài lấy
làm chuẩn thì nó mới chiếm một phần bảy. Câu điều kiện của bài đúng, và đúng
rộng rãi hơn bài tự nhận.

Số head: FLOP không đổi theo \(h\), đúng như bài nói, và đồng hồ thì đổi 5.37
lần từ \(h = 1\) tới \(h = 32\) ở \(n = 512\). Món này trả xong, kèm một đính chính cho chương
trước.

Thứ tự layer norm: nửa lý thuyết đo được ở lúc khởi tạo, gradient của post-LN
lớn gấp khoảng năm lần và không nhỏ đi theo độ sâu. Nửa thực nghiệm thì không,
vì hai tầng không phải quy mô mà chuyện này xảy ra.

Warmup: hình dạng của \eqref{eq:ch08-warmup} dẫn ra được, và chỗ nó nối vào
gradient ban đầu thì rõ. Bảng bốn ô không tách được nó khỏi nhiễu seed, và mục
ấy nói thẳng ra.

\subsection*{Chỗ ba chương sau bắt đầu}

Ba khoảng trống mà kiến trúc này để lại là ba chương sau, và chương này chạm
được vào cả ba theo ba mức khác nhau. Đánh số chúng ở đây để phần còn lại của
mục gọi tên được: thứ nhất là chi phí bậc hai theo độ dài, thứ hai là không có
\tn{thiên kiến quy nạp}{inductive bias} về vị trí, thứ ba là đói dữ liệu.

\index{chi phí!cái đắt nhất không phải FLOP}%
Khoảng trống thứ nhất là khoảng trống chương này đo, và phát hiện tôi không chờ
đợi khi bắt đầu nằm ở mục~\ref{sec:ch08-cho-vo}. Nó đổi cách tôi đọc cả bảng 1.

Chi phí thật của phần bậc hai, trên máy, không đến khi số phép tính vượt lên. Nó
đến khi \tn{ma trận điểm số}{score matrix} không còn nằm vừa trong cache, và
ngưỡng ấy tính bằng
byte chứ không tính bằng \(n\). Giữ nguyên \(n\) mà tăng batch thì vẫn gặp đúng
cái dốc ấy, ở đúng số byte ấy. Bảng 1 có ba cột và không cột nào đo cái đó, vì
cả ba cột đều là hàm của \(n\), \(d\), \(k\) và \(r\), và không cột nào biết máy
có bao nhiêu cache.

Mục~\ref{sec:ch08-cho-vo} đã nói FlashAttention thắng vì nó sửa đúng vào đại
lượng ấy. Chỗ này chỉ thêm một câu: đó là lý do phép đo nằm ở chương này chứ
không ở chương~\ref{ch:chi-phi-bac-hai}. Một chương mở đầu bằng \enquote{cách
này nhanh hơn} thì đang hứa; một chương mở đầu bằng một cột số không phẳng thì
đang nợ người đọc một lời giải thích, và đó là thứ dễ trả hơn nhiều.

\index{chi phí!thiên kiến quy nạp vẫn chưa trả}%
Khoảng trống thứ hai thì chương này không đụng tới, và điều nó đóng góp được là
một phép loại trừ. Mục~\ref{sec:ch07-nhieu-dau} và mục~\ref{sec:ch07-corpus} để
lại một nghi vấn hợp lý: nếu Transformer cần gấp ba số epoch mới đuổi kịp, biết
đâu cái nó thiếu chỉ là ngân sách tính toán, và bảng chi phí của chương này sẽ
nói ra điều đó.

Nó không nói. Bốn mục đầu đo chi phí theo ba cách và không cách nào tìm ra chỗ
kiến trúc bị phạt oan: ở độ dài câu mà corpus này chạy, phần bậc hai chưa tới
một phần trăm một tầng, ma trận điểm số nằm gọn trong cache, và số bước là số
bước. Cái thiếu không phải chỗ chứa cũng không phải thời gian mỗi bước. Loại trừ
xong thì chỉ còn lại thứ mà bạn không mua được bằng đồng hồ, và
chương~\ref{ch:bias-quy-nap} là chương dành riêng cho nó.

\index{chi phí!chỗ dữ liệu mua được thiên kiến}%
Khoảng trống thứ ba thì chương này chỉ chạm được vào mép, và mép ấy là
mục~\ref{sec:ch08-dieu-chuan}. Dropout làm hỏng ở đây và đáng 1.2 BLEU ở bài.
Cùng một cơ chế, cùng một giá trị tham số, hai kết luận ngược nhau, và cái phân
biệt chúng là 6000 câu sinh từ văn phạm so với 4.5 triệu cặp câu thật.

Đó là lần đầu trong sách này một kết luận lật hẳn chiều chỉ vì lượng dữ liệu, và
nó là bản xem trước của luận điểm trung tâm ở chương~\ref{ch:tien-huan-luyen} và
chương~\ref{ch:vit}. Ở chương~\ref{ch:vit} câu chuyện ấy đạt dạng gọn nhất của
nó: cùng một kiến trúc thua ResNet trên ImageNet-1k và thắng sau khi được huấn
luyện trước trên JFT-300M~\autocite{dosovitskiy2021vit}. Không đổi một dòng mã
nào.

Ba khoảng trống ấy là ba chương, và chương này đo được khoảng trống thứ nhất rồi
đo tiếp rằng nó không phải khoảng trống đắt nhất.
